\documentclass{sprz}

\addbibresource{bibliography.bib}

\studfield{Informatyka}
\studtype{Zaoczne}
\title{Praca inżynierska o tym i owym}
\engtitle{Something about this and that}
\acronym{BRR}
\titledate{2023-10-10}
\supervisor{dr Puźniakowski Tadeusz}
\reviewer{prof. dr hab. inż. lek. Jamroży Omnibus}
\author{Nowak Jan}{s1234}{Sztuczna Inteligencja}{Stacjonarny}
\author{Nowak Jan}{s1234}{Sztuczna Inteligencja}{Stacjonarny}
\author{Nowak Jan}{s1234}{Sztuczna Inteligencja}{Stacjonarny}
\author{Nowak Jan}{s1234}{Sztuczna Inteligencja}{Stacjonarny}
\author{Nowak Jan}{s1234}{Sztuczna Inteligencja}{Stacjonarny}
\consultant{--- brak ---} % Koniecznie trzeba podać brak, albo wpisać konsultantów tak jak przy autorach
\projectgoals{Celem tego projektu jest pokazanie jak można użyć klasy sprz do przygotowania pracy inżynierskiej na PJATK}
\productsandservices{Tu wpisujemy oczekiwane produkty}
\mainfunctionalities{Tu wpisujemy oczekiwane funkcjonalności}
\successmeasure{Kompilujący się elegancko szablon w postaci stylu LaTeX-owego który można użyć do robienia pracy inż.}
\projlimitations{Wymagana podstawowa znajomość systemu składu tekstu LaTeX.}
\date{miesiąc, rok obrony}
\finishdate{\today}
\nabstract{To jest streszczenie niezwykle innowacyjnej pracy inżynierskiej w której przeplatają się elementy paplaniny z wygenerowanym tekstem pseudonaukowym rozpoznawanym jako dezinformacja.}
\keyword{grafika komputerowa}
\keyword{komputer}
\keyword{XXI wiek}
\keyword{skład komputerowy}
\keyword{grafika 3D}



\begin{document}

\maketitle

\makeprojectcard

\makedeclarationsai

\tableofcontents

\chapter{Informacje wstępne}
\section{Cel i zakres pracy}
\section{Ogólny opis rozwiązania}
\section{Struktura pracy}

\chapter{Słownik pojęć}

\chapter{Opis problemu}
\section{Charakterystyka rynku wydarzeń motorsportowych}
\section{Problemy związane z organizacją wydarzeń i zapisami uczestników}
\section{Problemy informacyjne po stronie kierowców i organizatorów}
\section{Uzasadnienie potrzeby stworzenia systemu RacePortal}
\section{Aspekty techniczne i organizacyjne problemu}

\chapter{Analiza SWOT}
\section{Mocne strony (Strengths)}
\section{Słabe strony (Weaknesses)}
\section{Szanse (Opportunities)}
\section{Zagrożenia (Threats)}

\chapter{Kontekst biznesowy}
\section{Segmentacja użytkowników}
\section{Interesariusze systemu}
\section{Propozycja wartości systemu}
\section{Kanały komunikacji i dystrybucji}
\section{Kluczowe zasoby i działania}
\section{Model działania platformy}
\section{Realistyczność wdrożenia i dalsze perspektywy}
\section{Kontekst prawny i etyczny}

\chapter{Planowanie projektu}
\section{Cele projektu}
\section{Zakres projektu}
\section{Metodyka pracy}
\section{Organizacja zespołu}
\section{Harmonogram pracy}
\section{Infrastruktura i środowisko pracy}
\section{Analiza ryzyka projektowego}
\subsection{Analiza zagrożeń}

\chapter{Analiza wymagań}
\section{Źródła i sposób pozyskania wymagań}
\section{Wymagania ogólne i dziedzinowe}
\section{Wymagania funkcjonalne}
\section{Wymagania niefunkcjonalne}
\section{Wymagania bezpieczeństwa}
\section{Wymagania interfejsu użytkownika}
\section{Diagramy UML}
\subsection{Diagram przypadków użycia}
\subsection{Diagram klas}
\subsection{Diagram sekwencji: rejestracja i logowanie}
\subsection{Diagram sekwencji: zapis na wydarzenie}
\subsection{Diagram sekwencji: zarządzanie garażem}

\chapter{Projekt systemu}
\section{Założenia architektoniczne}
\section{Architektura logiczna systemu}
\section{Architektura fizyczna i wdrożeniowa}
\section{Projekt bazy danych}
\section{Projekt interfejsu REST API}
\section{Projekt warstwy webowej}
\section{Projekt aplikacji mobilnej}
\section{Projekt mechanizmów bezpieczeństwa}

\chapter{Realizacja}
\section{Organizacja implementacji}
\section{Podział systemu na komponenty}
\subsection{Backend}
\subsection{Frontend webowy}
\subsection{Aplikacja mobilna}
\section{Integracja komponentów}
\section{Konteneryzacja i uruchamianie środowiska}

\chapter{Szczegóły implementacji backendu}
\section{Struktura projektu backendowego}
\section{Model domenowy i encje}
\section{Rejestracja użytkowników i logowanie}
\section{Zarządzanie profilami użytkowników}
\section{Obsługa wydarzeń}
\section{Obsługa zgłoszeń}
\section{Obsługa garażu}
\section{Panel administratora i organizatora}
\section{Integracja z bazą danych}
\section{Obsługa wyjątków i walidacja}

\chapter{Szczegóły implementacji frontendu webowego}
\section{Struktura aplikacji webowej}
\section{Strona główna i katalog wydarzeń}
\section{Widok szczegółów wydarzenia}
\section{Logowanie i rejestracja}
\section{Dashboard użytkownika}
\section{Widok garażu}
\section{Panel organizatora}
\section{Panel administratora}
\section{Mapa i elementy lokalizacyjne}
\section{Responsywność i doświadczenie użytkownika}

\chapter{Szczegóły implementacji aplikacji mobilnej}
\section{Założenia aplikacji mobilnej}
\section{Architektura rozwiązania mobilnego}
\section{Lista wydarzeń i widok szczegółów}
\section{Nawigacja i interfejs użytkownika}
\section{Integracja z API}
\section{Przekierowanie do wersji webowej}
\section{Ograniczenia wersji mobilnej}

\chapter{Bezpieczeństwo systemu}
\section{Założenia bezpieczeństwa}
\section{Klasyfikacja danych}
\section{Uwierzytelnianie użytkowników}
\section{Autoryzacja i kontrola dostępu RBAC}
\section{Wykorzystanie tokenów JWT}
\section{Walidacja danych wejściowych}
\section{Ochrona zasobów użytkowników}
\section{Bezpieczeństwo komunikacji}
\section{Analiza zagrożeń}
\section{Analiza ryzyka}
\section{Ograniczenia przyjętych zabezpieczeń}

\chapter{Testowanie i weryfikacja systemu}
\section{Strategia testowania}
\section{Środowisko testowe}
\section{Testy jednostkowe}
\section{Testy integracyjne backendu}
\section{Testy End-to-End aplikacji webowej}
\section{Testy aplikacji mobilnej}
\section{Testy mechanizmów bezpieczeństwa}
\section{Automatyzacja uruchamiania testów}
\section{Analiza wyników testów}
\section{Ograniczenia procesu testowania}

\chapter{Prezentacja systemu}
\section{Prezentacja głównych funkcjonalności}
\section{Scenariusz użycia z perspektywy kierowcy}
\section{Scenariusz użycia z perspektywy organizatora}
\section{Scenariusz użycia z perspektywy administratora}
\section{Podsumowanie działania systemu}

\chapter{Nakład pracy}
\section{Liczba godzin przeznaczona na realizację projektu}
\section{Podział zadań w zespole}
\section{Narzędzia wspierające zarządzanie pracą}
\section{Problemy organizacyjne i ich rozwiązanie}

\chapter{Podsumowanie}
\section{Największe wyzwania i sposób ich rozwiązania}
\section{Zmiany i adaptacje w trakcie realizacji}
\section{Ograniczenia końcowej wersji systemu}
\section{Możliwości dalszego rozwoju}
\section{Wnioski końcowe}

\printbibliography[title={Bibliografia}, heading=bibintoc]

\makethesisattachments







\end{document}
